\subsection{RQ1: General generation approaches}
\label{sec:rq1}
\textbf{Statistical models} sample fitted distributions and dependencies,
for example through \glspl{bayesian-network}. Their usefulness depends on
whether those assumptions retain task-relevant relationships
\parencite{hernandezSyntheticDataGeneration2022,hansenReimaginingSyntheticTabular2023}.
More samples from a fitted generator do not automatically remove uncertainty
in that generator \parencite{breugelSyntheticDataReal2023}.

\textbf{Procedural and physical simulation} constructs scenes from rules
and assets, then models their observation. It offers scenario control and
access to labels, but omitted assets or sensor effects remain limitations.
\Gls{domain-randomization} varies parameters to broaden training conditions;
the chosen sampling distribution determines whether that variation is useful
\parencite{nikolenkoSyntheticDataDeep2019}.

\textbf{Learned generators} include \glspl{vae}, \glspl{gan} and
\glspl{diffusion-model}. They learn structure from examples, with different
latent, adversarial or denoising objectives. Conditional generation exposes
selected controls, while coverage and \gls{mode-collapse} remain concerns
\parencite{hernandezSyntheticDataGeneration2022,hansenReimaginingSyntheticTabular2023}.
LiDARGen uses \gls{score-matching} and iterative
\gls{langevin-dynamics}; its representation and sampling schedule are
therefore practical design choices
\parencite{zyrianovLearningGenerateRealistic2022}.

\textbf{Hybrids} combine explicit structure with learned corrections.
LiDARsim, for example, couples reconstructed assets and \gls{ray-casting}
with learned \gls{ray-drop}
\parencite{manivasagamLiDARsimRealisticLiDAR2020}.
RQ1 thus identifies a trade-off between explicit control and learned
regularities, qualified by coverage, annotation support, real-data dependence
and computation. The evidence does not establish a universally best family.
